\documentclass[10pt,a4paper]{amsart}
\usepackage[margin=19mm]{geometry}
\usepackage{amsmath,amssymb,mathtools}
\usepackage[scr=rsfs,cal=euler]{mathalfa}
\usepackage{xfrac}
\usepackage[unicode,hidelinks,bookmarks=false]{hyperref}
\newcommand{\ie}{i.e.\ }
\newcommand{\cf}{cf.\ }
\newcommand{\resp}{respectively\ }
\newcommand{\Subsection}{\subsection}
\providecommand{\nobreakdash}{\nobreak\hbox{-}}
\setlength{\emergencystretch}{2em}
\multlinegap0pt
\newcommand{\Cat}{\mathcal{C}}
\newcommand{\Der}{\mathcal{D}}
\newcommand{\End}{\operatorname{End}}
\newcommand{\HOM}{\operatorname{HOM}}
\newcommand{\Hom}{\operatorname{Hom}}
\newcommand{\kk}{\Bbbk}
\def\udgmod{\mbox{-\underline{gmod}}}
\theoremstyle{definition}
\newtheorem{thm}{Theorem}[section]
\newtheorem{cor}[thm]{Corollary}
\newtheorem{conj}[thm]{Conjecture}
\newtheorem{lem}[thm]{Lemma}
\newtheorem{rem}[thm]{Remark}
\newtheorem{prop}[thm]{Proposition}
\newtheorem{claim}[thm]{Claim}
\newtheorem{defn}[thm]{Definition}
\newtheorem{example}[thm]{Example}
\newtheorem{question}[thm]{Question}
\newtheorem{conv}[thm]{Convention}
\newtheorem*{thm*}{Theorem}
\title[Stable Cellularity in Hopfological Algebra]{Stable Cellularity in Hopfological Algebra}
\author{You Qi}
\date{Source: arXiv:2609.28295v1 (CC BY 4.0)}
\begin{document}
\maketitle

\noindent\emph{Source and licence.} The delivered work is the article shown in the title above, version arXiv:2609.28295v1 (CC BY 4.0), distributed under the Creative Commons Attribution 4.0 International licence, as recorded in the arXiv licence metadata for this exact version and shown on the abstract page saved with this item. The full licence notice, which carries the licence URL and the address of that abstract page, is the bundled file \texttt{licenses/arxiv-2609.28295-CC-BY-4.0.txt}.
\par\smallskip

\noindent\textit{Editorial scope.} Counted unit: the article's Lemma (source label lem:A-degree-general) (source0.tex lines 1164--1172, source label \texttt{lem:A-degree-general}): ``For any , , and , the following conditions hold: (i) For every internal degree , ; (ii) ; (iii) is finite dimensional.''. It is delivered together with the article's complete original proof (source0.tex lines 1174--1244, one proof environment adjacent to the statement, 71 lines), copied line for line from the article's own text. Required context retained uncounted: lem:Hom (lines 559--576); lem:positive-syzygies (lines 1121--1127). Editorial formatting only: an AMS \texttt{amsart} preamble that restates the article's own macro definitions and its own theorem declarations, and the theorem environments the retained lines use; the article's own class file is neither shipped nor input; the retained environments are renumbered sequentially in order of appearance, whereas the article prints them with its own numbering; the article's equation numbering is not reproduced and every cross-reference inside the retained text resolves to the wrapper's numbering. No citation occurs inside any retained line, so the wrapper carries no bibliography; All retained bodies are exact copies of the bundled \texttt{sources/211571/source0.tex}, so no omission receipt applies. No asset-inclusion command occurs inside any retained span and the package ships no image asset.


\section*{Retained context: lem:Hom (source0.tex lines 559--576)}

\begin{lem}\label{lem:Hom}
Let $M$ be $A$-projective and $N$ a graded $B$-module.  The graded space
$\HOM_A(M,N)$ carries the standard $H$-module structure, and
\[
\Hom_B^r(M,N)=\Hom_H^r(\kk,\HOM_A(M,N))= \Hom_A^r(M, N)^H .
\]
If $M$ is cofibrant, then morphisms from $M$ in the homotopy and derived
categories agree, and for every internal degree $r$ one has
\begin{equation}\label{eq:HOM-formula}
\Hom_{\Der(A,H)}^{r}(M,N)
=
\Hom_{\Cat(A,H)}^{r}(M,N)
\cong
\dfrac{\Hom_B^r(M,N)}
{\Lambda\cdot \Hom_A^{r-\ell}(M,N)} .
\end{equation}
Here $\Lambda$ is a homogeneous integral of degree $\ell$.  
\end{lem}


\section*{Retained context: lem:positive-syzygies (source0.tex lines 1121--1127)}

\begin{lem}\label{lem:positive-syzygies}
For $0\le i<\ell$ and $j>0$, the graded $H$-module $V_i[-j]$ is isomorphic to an object supported in strictly positive degrees in $H\udgmod$.  Consequently
\begin{equation}\label{eq:H-positive-bound}
\Hom_{H\udgmod}^{r}(V_a,V_b[-j])=0
\qquad(j>0,\ r\le0).
\end{equation}
\end{lem}


\section{The counted unit: Lemma (source label lem:A-degree-general) and its complete original proof}

\begin{lem}\label{lem:A-degree-general}
For any $x,y\in X$, $0\le a,b<\ell$, and $j>0$, the following conditions hold:
\begin{enumerate}
\item[(i)] For every internal degree $r\le0$,
$\Hom_{\Der(A,H)}^r(E_{x,a},E_{y,b}[-j])=0$;
\item[(ii)] $\End_{\Der(A,H)}(E_{x,a})=\kk$;
\item[(iii)] $\Hom_{\Der(A,H)}(E_{x,a},E_{y,b})$ is finite dimensional.
\end{enumerate}
\end{lem}

\begin{proof}
For (i), choose a finite-dimensional graded $H$-module $W$, supported in
strictly positive degrees, representing $V_b[-j]$ in $H\udgmod$, as in
Lemma~\ref{lem:positive-syzygies}.  Then
$E_{y,b}[-j]\cong P_y\otimes W$ in $\Der(A,H)$.  Since
$P_x\otimes V_a$ and $P_y\otimes W$ are cofibrant, Lemma~\ref{lem:Hom} gives
\begin{equation}\label{eq:Hom-as-stableHom}
\Hom^{r}_{\Der(A,H)}(E_{x,a},E_{y,b}[-j])
\cong
\dfrac{\Hom^r_B(P_x\otimes V_a,P_y\otimes W)}
{\Lambda\cdot
\Hom_A^{r-\ell}(P_x\otimes V_a,P_y\otimes W)} .
\end{equation}
As graded vector spaces,
\begin{equation}\label{eq:A-Hom-degree}
\HOM_A(P_x\otimes V_a,P_y\otimes W)
\cong
e_xAe_y\otimes V_a^*\otimes W.
\end{equation}
The three tensor factors on the right are supported in degrees $\ge0$, $\ge0$, and
$>0$, respectively.  Hence the right-hand side of
\eqref{eq:A-Hom-degree} is supported in strictly positive degrees.  Therefore
\[
\Hom_A^r(P_x\otimes V_a,P_y\otimes W)=0
\qquad(r\le0),
\]
and in particular the numerator of \eqref{eq:Hom-as-stableHom} vanishes.
This proves (i).

For (ii), take $j=0$ and $r=0$ in Lemma~\ref{lem:Hom}.  The module $V_a$ is
cyclic, generated in its lowest degree $-a$, and
\[
(E_{x,a})_{-a}
=
A_0e_x\otimes (V_a)_{-a}
\cong\kk.
\]
Thus every degree-zero $B$-endomorphism of $E_{x,a}$ is a scalar multiple of
the identity:
\[
\End_B^0(E_{x,a})=\kk.
\]
Moreover,
\[
\HOM_A(E_{x,a},E_{x,a})
\cong e_xAe_x\otimes V_a^*\otimes V_a
\]
is supported in degrees $\ge-a$.  Since $a<\ell$, its degree $-\ell$ part is
zero.  Hence the denominator in \eqref{eq:HOM-formula} vanishes for $r=0$,
and
\[
\End_{\Der(A,H)}(E_{x,a})=\kk.
\]

For (iii), it is enough to note that
\[
\Hom_B^0(E_{x,a},E_{y,b})
\subseteq
\Hom_A^0(E_{x,a},E_{y,b}).
\]
Using
\[
\HOM_A(E_{x,a},E_{y,b})
\cong e_xAe_y\otimes V_a^*\otimes V_b,
\]
the degree-zero part involves only $A_d$ with $0\le d\le b$.  Since $A$ is
locally finite and $V_a,V_b$ are finite dimensional,
$\Hom_A^0(E_{x,a},E_{y,b})$ is finite dimensional.  The same is therefore
true after passing to the quotient defining
$\Hom_{\Der(A,H)}(E_{x,a},E_{y,b})$.
\end{proof}

\end{document}
